\documentclass[12pt]{letter}
\usepackage[margin=1in]{geometry}
\usepackage{amsmath}
\usepackage{hyperref}

\signature{Henry Laverde-Rojas \\ Carlos Laverde-Rodriguez}
\address{Programa de Econom\'ia \\ Universidad Militar Nueva Granada \\ Campus, Colombia}

\begin{document}

\begin{letter}{Editors-in-Chief \\ \emph{Decision Support Systems} \\ Elsevier}

\opening{Dear Editors,}

I am pleased to submit \emph{``The Post-ChatGPT Contraction
Reaches the Crowd-Validated Margin: A Within-Tag Triple
Difference and a Value-Calibrated Early-Warning System for
Stack Overflow''} for consideration at \emph{Decision Support
Systems}.

\noindent\textbf{Problem and decision relevance.}\\
Public question-and-answer communities such as Stack Overflow
are both a knowledge resource for millions of practitioners
and the training substrate for the very generative-AI systems
that now compete with them. Recent work establishes that
Q\&A activity contracted after the November 2022 ChatGPT
release (Burtch, Lee, and Chen 2024; del R\'io-Chanona,
Laurentsyeva, and Wachs 2024, $\sim$15--25\% declines) and
that the surviving question population is, on average, of
higher quality (Xue et al.\ 2026). A platform manager---or an
AI provider monitoring the health of its training
corpus---reading those two facts together might conclude that
generative AI prunes only low-value posts that public Q\&A
would not have retained as durable artefacts. If true, no
intervention is warranted. The decision-relevant question,
which an aggregate quantity-and-quality decomposition cannot
answer, is whether the contraction also reaches the
\emph{high-value, crowd-validated} posts that sustain the
platform's downstream knowledge function---and, if so, how a
manager who observes only activity counts should detect and
triage the loss. This paper answers both the measurement
question and the decision question.

\noindent\textbf{Contributions.}\\
\begin{enumerate}
\item \emph{The strict benign-pruning reading is rejected at
the within-tag value margin.} A single pre-specified shape
test over mutually exclusive net-score bins returns a
negative average displacement across the crowd-validated
range ($\hat\beta = -0.167$, $p<10^{-4}$) that is roughly
constant over the low-to-mid bins (score $0$, $1$, $2$:
$-0.21$, $-0.20$, $-0.21$) and attenuates only at the rare
elite tail (slope $p = 2\times10^{-4}$). The displacement is
therefore not confined to score-zero posts.
\item \emph{A within-tag, within-question-type triple
difference} at the (tag $\times$ week $\times$ question-type)
cell isolates the LLM-substitutability slice of the
aggregate decline, reported under \emph{two co-primary}
treatment measures---a historical answerability index
($-0.108$, $p=0.021$) and an embedding-based score that uses
no historical features ($-0.119$, $p<0.001$)---so the result
cannot be dismissed as historical popularity relabelled.
\item \emph{A value-calibrated early-warning system---the
decision-support artefact.} We let the data choose the
detection signal: out of sample on the full 100-tag panel,
the activity-decline metric a platform already tracks
predicts realised 2025 high-value erosion at AUC $0.98$,
while a value-weighted detector reaches only $0.57$. The
system therefore \emph{detects} on activity, and the paper's
value-margin estimates serve as the \emph{calibration} that
converts an activity alert into the high-value posts at risk
that managers actually care about. The calibration earns its
place at the intensive (triage) margin: allocating
intervention effort by the calibrated high-value-at-risk
nearly attains the perfect-foresight optimum, whereas
allocating by raw activity leaves up to $5\%$ of the
protectable high-value flow unprotected. A highly persistent
composition flag ($\rho = 0.90$ from 2024 into the new 2025
data) classifies whether a tag's decline is
value-disproportionate, informing intervention \emph{type}.
\end{enumerate}

\noindent\textbf{Fit with \emph{Decision Support Systems}.}\\
The paper's core deliverable is a decision aid, evaluated as
one. The early-warning system is built from signals a
platform already collects, is calibrated on identified
causal estimates rather than correlations, and is assessed by
its out-of-sample detection performance and by a
value-of-information / regret analysis that quantifies the
cost of the naive activity-only policy a manager would
otherwise use. This is squarely the journal's remit:
analytics and AI in support of managerial decisions, with an
explicit account of the decision the tool improves and the
loss it averts. The empirical contribution (the within-tag
value-margin estimate) is what makes the calibration
trustworthy; the artefact is what makes it actionable.

\noindent\textbf{Identification, stated honestly.}\\
Formal pre-trend tests reject \emph{exact} parallel trends,
so we do not claim a standalone natural experiment: the
estimates are conditional post-ChatGPT associations aligned
with the LLM-substitutability mechanism. Their credibility
rests on a conjunction of diagnostics that fail under
different threats---a sign reversal between the real cutoff
and 29 rolling pre-period placebos, a positive (placebo)
June 2022 GitHub Copilot cutoff, monotone growth of the
event-study coefficient with no reversion, and a HonestDiD
$\Delta^{\text{RM}}$ sensitivity bound under the full
event-study variance--covariance matrix with a breakdown at
$\bar M = 0.75$, which we read as a \emph{moderate} signal
and say so. The estimate also survives a 32-cell
pre-specified specification curve (all negative, all
significant), PPML, two-way clustering, and a wild cluster
bootstrap. Because two one-off 2023 platform events (a
data-licensing change and a moderator strike) overlap the
early post-period, we treat a donut difference-in-differences
that excludes the affected weeks as the \emph{primary} test
against that confound ($-0.110$, $p=0.022$). Finally,
re-estimating on the \emph{full 100-tag panel extended
through 2025} strengthens the triple difference to $-0.172$
($p=0.002$) and the quarterly association deepens to $-0.68$
by the end of 2025 with no reversion, removing the reading
that the effect reflects a transient late-2024 episode.

\noindent\textbf{Scope discipline.}\\
We resist overclaiming. The identified channel is bounded at
$\sim$4\% of post-period activity (and $\sim$5\% of a
trend-implied shortfall); a five-transition support-post
funnel translates it into a cumulative shortfall of
$\sim$28{,}000 accepted, non-closed, non-negative-score posts
(about 40\% of the activity shortfall). The estimate is an
innovation \emph{input} rather than an outcome, and the paper
states plainly what the magnitude does \emph{not} support
(welfare, innovation outcomes, or generalisation beyond
programming Q\&A). A conditional-resolvability null under
fine-grained fixed effects shows that the descriptively
documented quality improvement is consistent with selection
into the survivor population rather than within-cell
upgrading. All references have been verified against
Crossref/DOI.

\noindent\textbf{Reproducibility and structure.}\\
The full collection-and-analysis pipeline is reproducible
from public, zero-budget data: manual Stack Exchange Data
Explorer exports through May 2023 and the public Stack
Exchange API v2.3 thereafter (including the 2025 extension),
with byte-level validation across the source migration on the
overlap period. The manuscript is self-contained within the
journal's 34-page limit (double-spaced, 1-inch margins); to keep
it within that limit the full robustness battery, the per-bin
value-distribution detail, the mechanism and
classifier-validation detail, the descriptive counterfactual,
and the literature-positioning table are provided in the public
replication package rather than as journal-hosted supplementary
material. All Python scripts, panels, tables, figures, and the
\LaTeX{} source are documented in
\texttt{docs/replication\_notes.md} and available at
\url{https://github.com/hlaverde/so-genai-value-margin}.

\noindent\textbf{Suggested expertise.}\\
Referees with expertise in decision-support and early-warning
systems, analytics for platform governance, the economics of
generative AI, online knowledge communities, or
quasi-experimental methods (triple difference, pre-trend
sensitivity bounds, cluster-robust inference, specification
curve analysis) would be well positioned to evaluate the
work.

\noindent\textbf{Disclosures.}\\
The work is single-authored and has not been submitted to any
other journal. There are no conflicts of interest. No paid
APIs, restricted institutional data, or external funding were
used.

\closing{Sincerely,}

\end{letter}
\end{document}
